% SIMPAT Sections 1-2 (P5 S1, 12-10-2026): simulation-credibility framing. Results 1-4 and all numbers unchanged
% from ress_sec12.tex (v1.2); only framing, related-work groups and references changed.
\section{Introduction}\label{sec:intro}

Simulation models are used more and more as evidence that a system may be accepted; regulators, for instance, have
started to accept evidence obtained in silico~\cite{viceconti2021insilico}. When such evidence is credible has its own
literature. Validation compares a model with real data for its intended use~\cite{sargent2013vv,oberkampf2010vv,oberkampf2004vvcomp};
the ASME V\&V~40 standard assesses the credibility a model needs through a risk-based framework~\cite{asme2018vv40,aldieri2023vv40};
verification, validation and uncertainty quantification frameworks and Bayesian calibration attach uncertainty to model
predictions~\cite{roy2011vvuq,kennedy2001calibration}; and the credibility of digital twins has been evaluated with
evidence theory~\cite{shi2025dtcredibility}. These frameworks express credibility qualitatively, as levels or scores,
or through uncertainty analysis. None of them states how many field observations a credible model may replace in a
given acceptance test.

A perception system on a fixed camera is a favourable case for this question: the context of use is the acceptance of a new
fixed-camera installation within a declared domain of static scenes, and a composite model can keep the real background, illumination and detector while synthesising only
the passing object, so operational validation of the composite model, for this decision, reduces to comparing the system's output on
paired real and synthetic events. We
study the question for one common test. A zero-failure (success-run) test accepts a unit when it passes a planned
number of demands (independent trials; here, object passages) without failure~\cite{hanley1983,meeker2017}. If each demand fails independently with probability $p$,
a unit with $p>\varepsilon$ passes $n$ demands with probability at most $(1-\varepsilon)^n$, so passing
$n_0(\varepsilon,\delta)=\lceil\ln\delta/\ln(1-\varepsilon)\rceil$ demands without failure demonstrates $p\le\varepsilon$
at confidence $1-\delta$: \nNzero{} demands at $\varepsilon=\delta=5\%$. In reliability engineering this test is called
a reliability demonstration test; we call it a \emph{zero-failure acceptance test}. Its sample-size burden motivates
test plans that borrow auxiliary information~\cite{ding2024truncated,wang2024monoprop}.

The unit under test is a learned perception component: a fixed object detector on a fixed camera that samples frames
periodically. A demand is the passage of an object (an \emph{event}); a failure is an event on which no sampled frame
yields a detection. Demands are rare and every new installation is a new unit, so \nNzero{} failure-free events must be
collected per camera; for road vehicles the field exposure needed to show rare-failure reliability is
prohibitive~\cite{kalra2016}. Simulation is the obvious remedy: digital twins already serve as platforms for testing
path-planning algorithms~\cite{wang2025dtplatform}, as virtual scenes for camera detectors~\cite{li2025flamedt}, and as
generators of synthetic sensor data with a validation step~\cite{nosenko2026gammadt}; virtual scenario
testing is a leading proposed alternative to mileage-based testing of automated vehicles~\cite{riedmaier2020}. For
perception, however, the simulation-to-reality gap is measurable and depends on conditions~\cite{reway2020simgap,dieter2023drone},
and performance in simulation need not predict performance in the field~\cite{kadian2020}. The credibility question
is therefore quantitative: \emph{how many of the $n_0$ field events can simulation evidence replace, and what form of
guarantee survives when the evidence is transferred to an untested unit?}

This paper answers the question for the zero-failure acceptance test of a perception component. We use the
\emph{composite simulation model}, or \emph{synthetic twin}, of an event: the real background of the event with the
object replaced by a rendered one drawn on the ground-truth trajectory (Section~\ref{sec:protocol}).
\begin{enumerate}
\item \emph{No free credit.} Without an assumption that links the simulator to reality, no valid procedure can use
independent simulated data to shorten a zero-failure test, for any simulator and any amount of simulation.
This makes precise, for success-run tests, a principle known from clinical
borrowing~\cite{koppschneider2020}, and relates to the limits on combining evidence in software
dependability~\cite{littlewood1993}.
\item \emph{The exact price of relaxing confidence.} Accepting error probability $\delta'>\delta$ when the simulator is
wrong buys at most, and exactly, $n_0(\varepsilon,\delta)-n_0(\varepsilon,\delta')$ demands: \ncExactA, \ncExactB,
\ncExactC{} and \ncExactD{} of the \nNzero{} for $\delta'=7.5$, 10, 15 and 20\%.
\item \emph{A paired-twin bound and a fleet-level transfer bound.} A synthetic twin of each real event bounds the real
failure probability by the twin failure probability plus a discordance rate, with no assumption; on one camera this
costs more real events than the plain test (\nCthreeA--\nCthreeC). Calibrated once across scenes, the discordance
transfers to an untested camera as a population-average bound whose price is the number of calibration scenes, not the
number of events.
\item \emph{An exchange rate on public benchmarks, and a named failure mode.} On \nMainScenes{} fixed-camera scenes the twin
as built saves no real event for the median camera (\nTcamfinOPS{} of \nMainScenes{} cameras save any); a twin run to convergence would save at most \nSavedOP{} real events per camera, if the discordance calibrated with
paired twins also held for the twins of new cameras (not verified), at the
operating point fixed in advance, where real failures are rare (\nPmainOP) and the certificate is therefore untested.
The transfer certificate stands at scene level but is withdrawn when the unit is the site (\nSiteSm{} sites).
When out-of-domain cameras are judged alone, they pass through: \nNightFC{} of \nNightScenes{} night cameras and, at a
fixed sampling phase, \nOODSJitfa{} of \nJitterScenes{} camera-jitter cameras (a scene with real shake) are falsely accepted; we call this
\emph{out-of-domain pass-through}, with two named modes, night and camera jitter.
\end{enumerate}

\emph{Simulation credibility is today scored qualitatively or by uncertainty analysis; to our knowledge, no result states how much field
evidence a credible twin may replace in a zero-failure acceptance test, nor how that credit transfers to an untested
unit. This paper gives a finite-sample answer (none without an assumption; an exact price for relaxing confidence; a
fleet-level bound whose price is scene count) and measures it on \nMainScenes{} fixed-camera scenes.}

% =====================================================================================================
\section{Related work and research gap}\label{sec:related}

Table~\ref{tab:gap} positions the paper against four bodies of work.

% Table 1 (SIMPAT): research gap, five rows. Row (e) numbers via numbers.tex (same macros as tab_gap.tex).
\begin{table*}[t]
\centering
\caption{Research gap. Credibility: how the worth of simulation evidence is expressed. ZF: is the zero-failure acceptance
test treated? Field events: is the number of field observations that simulation replaces quantified? Transfer: does the
guarantee reach an untested unit? Link: assumption needed between simulation and reality. Evidence: empirical basis.}
\label{tab:gap}
\scriptsize
\setlength{\tabcolsep}{2pt}
\begin{tabular}{@{}L{0.17\textwidth}L{0.13\textwidth}L{0.04\textwidth}L{0.15\textwidth}L{0.15\textwidth}L{0.15\textwidth}L{0.13\textwidth}@{}}
\toprule
Body of work & Credibility expressed as & ZF & Field events replaced & Transfer & Sim--real link & Evidence \\
\midrule
(a) Simulation validation and credibility~\cite{sargent2013vv,oberkampf2010vv,asme2018vv40,kennedy2001calibration,roy2011vvuq,pathmanathan2017applicability,shi2025dtcredibility}
 & level, score or predictive uncertainty & no & not quantified & applicability to a new context of use argued, not bounded & validation against field data; discrepancy model & case studies \\
(b) Simulation-assisted testing of automated and perception systems~\cite{wang2025dtplatform,li2025flamedt,riedmaier2021vvuq,reway2020simgap,luo2025sim2val,angelopoulos2023ppi}
 & gap measured; variance reduction for means & no & fewer real samples for a mean metric & not addressed (extrapolation named as open) & paired outputs (control variate, rectifier) or none & benchmarks, case studies \\
(c) Zero-failure test planning~\cite{hanley1983,baze1979,wilson2021assurance,ding2024truncated,wang2024monoprop,zhao2018perfection}
 & prior or auxiliary information & yes & through a prior; not as field events saved & to the tested unit; prior-system evidence handled conservatively & prior or similarity assumption (Bayesian); none (frequentist) & case studies \\
(d) Finite-sample guarantees and assurance for learned perception~\cite{angelopoulos2025ltt,degrancey2022conformaldet,mei2025pwc,yuan2026conformal,paterson2025amlas}
 & not addressed (real calibration data) & no & none & exchangeability of calibration and deployment & exchangeability & detection data sets \\
(e) This paper & field events saved per unit & yes & exact: none without a link; $n_0(\delta)-n_0(\delta')$ when relaxing; \nSavedOP{} events per camera at most & yes at scene level (withdrawn at site level), priced by scene count ($\UCP(J,m)=\nMainUpop$, $J=\nMainJ$, $m=\nMainScenes$); scenes exchangeable (i.i.d.), as in (d) at scene level & paired twin + measured discordance $D$ & \nMainScenes{} fixed-camera scenes, \nMainN{} events; failure mode measured \\
\bottomrule
\end{tabular}
\end{table*}


\emph{(a) Simulation validation and credibility.} Verification and validation of simulation models is an established
discipline that separates conceptual model validity, verification, operational validity and data
validity~\cite{sargent2013vv,oberkampf2010vv}. In validation the relation between computation and experimental data is
the issue, and predictive capability goes beyond it~\cite{oberkampf2004vvcomp,roy2011vvuq}; Bayesian calibration
learns the unknown parameters of a computer model from physical observations of the system~\cite{kennedy2001calibration}.
ASME V\&V~40 provides a risk-based framework for assessing the credibility of a computational model through predefined
credibility activities~\cite{asme2018vv40,aldieri2023vv40}, and a recurring difficulty is that validation cannot be
performed in a setting that closely resembles how the model will be used~\cite{pathmanathan2017applicability}. Regulatory
use of in silico evidence has made such credibility assessment a practical matter~\cite{viceconti2021insilico}. For digital
twins, credibility has been evaluated with evidence theory~\cite{shi2025dtcredibility}; twins are verified against
physical experiments~\cite{wang2025maglevdt,ye2023structuraldt}, driving-simulator validation studies tie the credibility
of results to scenario fidelity~\cite{zhang2025drivingsim}, and a review of manufacturing twins found little work that
performed both verification and validation~\cite{bitencourt2024dtvv}.
\emph{What is missing:} credibility is expressed as a level, a score or a predictive uncertainty; none of these is
converted into the number of field observations a credible model may replace in a given acceptance test, and none
bounds how that replacement carries over to a unit whose own validation data are absent.

\emph{(b) Simulation-assisted testing of automated and perception systems.} Digital twins serve as platforms for
testing and optimising path-planning algorithms~\cite{wang2025dtplatform}, as virtual tunnels for a camera flame
detector~\cite{li2025flamedt} and as generators of synthetic spectra with a validation step~\cite{nosenko2026gammadt};
a multifidelity simulation audit screens the dynamic executability of continuum-robot motions and reports
accepted cases out of randomized trials~\cite{wang2026simaudit}; virtual scenario
testing has been proposed in place of mileage-based testing of automated vehicles~\cite{riedmaier2020}, and model
validation for automated driving names extrapolation to untested conditions as a major challenge~\cite{riedmaier2021vvuq}.
The simulation-to-reality gap of camera object detectors has been measured across day, night, fog and
rain~\cite{reway2020simgap} and for drone detection~\cite{dieter2023drone}; synthetic images used to test driving
perception have a distribution mismatch that can bias test results~\cite{amini2024translators}, perception digital twins
reproduce field data only approximately~\cite{saad2025dtperception}, statistical model checking evaluates perception in
simulation only~\cite{barbier2019smc}, simulation performance can fail to predict field performance~\cite{kadian2020},
and a statistical foundation for aligning synthetic and real testing outcomes has been called
for~\cite{zhao2025statfoundation}. Closest to our question, Sim2Val uses simulator outputs as control variates to reduce
the number of real samples needed for a confidence bound on mean performance (CoRL 2025, arXiv
preprint)~\cite{luo2025sim2val}, SCAPE combines limited real rollouts with simulation proxies for policy evaluation
(preprint)~\cite{zhu2026scape}, multi-fidelity fusion weights simulation by its fidelity gap in reliability
prediction~\cite{wang2026fusion}, and prediction-powered inference combines many proxy outcomes with few real labels
through a paired rectifier~\cite{angelopoulos2023ppi,fisch2024stratppi}. \emph{What is missing:} these methods estimate
a mean performance or reduce its variance; none treats the zero-failure regime, where the quantity to bound is a rare
failure probability and the rectifier itself would cost almost $n_0$ real events per unit.

\emph{(c) Zero-failure test planning.} The zero-numerator bound and the rule of three are
classical~\cite{hanley1983,meeker2017}, and the Bayesian zero-failure (BAZE) test sets the required test time through a
prior on the failure rate~\cite{baze1979}. Later work designs sample sizes by assurance~\cite{wilson2021assurance},
studies the limits of Bayesian zero-failure plans~\cite{yoon2025bayesian} and plans success--failure verification and
qualification tests with subsystem, Monte-Carlo or expert information~\cite{ding2024truncated,wang2024monoprop,tan2025weibull}.
For software and autonomous systems, conservative Bayesian inference uses prior knowledge, including evidence from
previous similar systems, without optimistic bias~\cite{zhao2018perfection,zhao2019cbi,zhao2020operational}, and reusing
the reliability of an older system as a prior is optimistic rather than conservative~\cite{littlewood2020replacement}.
Bayesian borrowing quantifies the worth of a prior through its effective sample size~\cite{morita2008ess} and discounts
historical data through power or commensurate priors~\cite{ibrahim2000power,hobbs2011commensurate}; conservative
frequentist bounds exist when the test profile differs from the operational profile~\cite{bishop2017profiles}.
\emph{What is missing:} a frequentist statement of how much the zero-failure sample size can be reduced by simulation
rather than prior-system evidence, and at what price.

\emph{(d) Finite-sample guarantees and assurance for learned perception.} Conformal prediction gives distribution-free
finite-sample guarantees~\cite{angelopoulos2023gentle}, calibrated risk control extends them to general
losses~\cite{angelopoulos2025ltt}, and they have been applied to object detection~\cite{degrancey2022conformaldet},
keypoint-based pose estimation~\cite{yang2023purse}, perception for safe navigation~\cite{mei2025pwc} and anomaly detection
in cyber-physical systems~\cite{yuan2026conformal}. Reliability of machine-learning components is also modelled directly,
through multi-version architectures~\cite{wen2025nversion}, safety-assurance cases for object
detectors~\cite{paterson2025amlas} and failure-free-operation definitions under an operational
profile~\cite{aghazadeh2026hipllm}. \emph{What is missing:} these guarantees are calibrated on real data from the
deployment distribution; they do not say how much of that calibration data may be synthetic, nor how a guarantee
transfers to a unit whose own data are absent.
